% Purpose: motivate control validation for radial-structure diagnostics and state the contribution.
\section{Introduction}
\label{sec:introduction}

Radial-structure diagnostics correlate a radially organized covariate with
the residuals of galaxy rotation-curve fits in order to probe inner-region
structure. Rotation curves and radial profiles are widely used observational
diagnostics in galaxy studies, including broad rotation-curve reviews,
high-resolution nearby-galaxy rotation-curve surveys, and galaxy-mass reviews
\citep{SofueRubin2001,deBlok2008,Courteau2014}. A practical difficulty is
that the radius itself, and any smooth radially organized covariate, can
induce apparent correlations that do not reflect localized structure. Establishing
whether a measured correlation
carries information beyond broad radial organization is therefore primarily
a computational and statistical problem of control design.

Permutation and surrogate procedures are established tools in general
statistical and correlation-testing contexts, with related resampling and
non-parametric methods widely discussed in astronomy-statistics practice
\citep{SchreiberSchmitz2000,Lancaster2018,FeigelsonBabu2012,PhipsonSmyth2010,YuHutson2022}.
They also have a lineage in surrogate-data testing of an observed statistic
against randomized alternatives \citep{Theiler1992}. These procedures are
general-purpose; we
do not claim that any particular control family is common published
practice in rotation-curve analysis. Different control families preserve
different structure, and a correlation that exceeds a label-shuffle
comparison may still be reproduced by controls that preserve radial
ordering.

{\setlength{\emergencystretch}{2em}
In this work we present a reproducible, branch-separated control-validation
workflow for radial-structure diagnostics in SPARC-like rotation-curve data
\citep{Lelli2016}. The implemented workflow is referred to here as the Nisa
Radial Control Battery (NRCB). The workflow evaluates a radius-conditioned
partial-rank statistic against stochastic shuffle/permutation empirical control
distributions and against exact circular-shift randomization distributions,
and evaluates an amplitude-based spectral-concentration statistic against
amplitude-preserving random-phase controls and fixed smooth references. Our
contribution is the design of this control battery, its deterministic-seed
implementation, and an audited descriptive demonstration on two explicitly
separated galaxy samples. We make no physical interpretation or
model-validation claim; the demonstration characterizes how the control
families behave, not the astrophysics of any galaxy.
\par}

The paper is organized as follows. Section~\ref{sec:data} defines the data
context and the two sample branches. Section~\ref{sec:methods} defines the
diagnostic statistics and control families. Section~\ref{sec:demonstration}
presents the descriptive demonstration. Section~\ref{sec:discussion} states
limitations, Section~\ref{sec:reproducibility} documents reproducibility,
and Section~\ref{sec:conclusion} concludes.



